\section{Positioning: what exists and the claimed gap}
\label{sec:positioning}

Most ingredients of this pipeline already exist in isolation, so the contribution cannot rest on
``soft preferences'', ``distributional DPO'', ``few human labels'' or ``LLM-judge calibration''.
Table~\ref{tab:positioning} lists the closest work.

\begin{table}[h]
\centering
\small
\renewcommand{\arraystretch}{1.15}
\begin{tabularx}{\linewidth}{@{}p{3.1cm}XX@{}}
\toprule
Work & What it establishes & What it leaves open for us \\
\midrule
RLAIF \newline \citep{lee2024rlaif} &
AI preference \emph{distributions} from normalized token log-probabilities (e.g.\ $[0.6,0.4]$);
reward model trained by cross-entropy on these soft labels. &
No few-shot correction toward humans; no guarantees. \\
GAPO \newline \citep{furuta2024gapo} &
DPO-style optimization for soft preference labels via a weighted geometric average of likelihoods. &
Takes soft labels as given; how accurate they are and where they come from is not modeled. \\
MixDPO \newline \citep{imai2026mixdpo} &
Latent distribution of preference \emph{strength} inside DPO. &
Heterogeneity of strength, not a judge-to-human correction. \\
DPRM \newline \citep{li2024dprm} &
Population preferences as a categorical distribution; OT-based reward model; policy fine-tuning. &
Not a few-shot correction of an AI judge; no rate theory. \\
RLTHF \newline \citep{xu2025rlthf} &
Targeted human corrections of AI-mislabeled samples; matches full human annotation with
6--7\% of the effort. &
Label-level correction; no distributional target, no rates. \\
\citet{vandenburg2025aligning} &
Linear map from black-box LLM outputs to human judgments, zero- and few-shot. &
Parametric; no minimax theory; no post-training. \\
\citet{shen2025active} &
Fisher-information selection of pairs for human annotation in reward modeling. &
Design is not tied to a judge-to-human discrepancy. \\
DIAL \newline \citep{cai2026dial} &
Joint BTL likelihood of abundant LLM and scarce human comparisons, weight $\lambda$ ($\lambda=0$ is
human-only) tuned by GACV; asymptotic normality. &
Binary verdicts; per-pair probabilities only through a constrained BTL parameterization of item
scores; parametric; no post-training. \\
\citet{galatolo2026lowresource} &
Linear probe on $\le 500$ human-labeled pairs annotates $50$K$+$ pairs for preference optimization. &
No judge distribution; empirical. \\
\citet{li2026personalizing} &
Black-box FSP: local smoothing, adaptive borrowing, sampling design, minimax and no-harm; includes
binary classification (so Bernoulli outputs). &
Scalar predictor as the offset; no graded judge outputs; Lebesgue-risk design; no post-training. \\
\bottomrule
\end{tabularx}
\caption{Closest prior work (checked 28 September 2026).}
\label{tab:positioning}
\end{table}

Two recent works are also relevant to individual components. \citet{li2026calibrate} calibrates
full panels of noisy judges with proper scoring rules; this is a natural calibration baseline for
Method~I. \citet{lail2026decomposing} targets expert labels using the epistemic part of judge
uncertainty; this is a baseline for Method~II.

\keybox{\textbf{Claimed gap.} Few-shot personalization of a \emph{graded, distribution-valued}
black-box judge, with minimax guarantees for \emph{bounded} preference targets and a population-risk
label design, carried through to an end-to-end minimax rate for the post-trained policy's excess
preference risk (Target~\ref{tr:e2e}).}

We did not find prior work that establishes this complete chain. We treat it as a
\emph{novelty hypothesis} that must be re-checked before submission. DIAL \citep{cai2026dial},
posted three days before this draft, is the closest on adaptive borrowing and should be
tracked.

\paragraph{Claims we do not make.}
The soft-label preference loss is not new \citep{lee2024rlaif,furuta2024gapo}. Neither is ordinal or
strength-aware preference optimization \citep{imai2026mixdpo}, active selection of pairs for human
labeling \citep{shen2025active,xu2025rlthf}, or mapping judge outputs to human labels
\citep{vandenburg2025aligning}. The nested-information identity of
Proposition~\ref{prop:scalar-vs-dist} is standard: its $G$ is a variable-importance numerator
\citep{williamson2021nonparametric,williamson2023general}, and the panel version is Proposition~1 of
\citet{li2026calibrate}. The bridge propositions of Section~\ref{sec:post-training} are standard
target-perturbation and strong-convexity arguments. Our contribution is where the preference target
comes from, the minimax and design theory for it, and what guarantees it carries through to the
policy.
